\documentclass[10pt,a4paper]{article}
\usepackage{cvstyle}

\graphicspath{{./}{src/}{src/apps/}{plantilla/PIC/}}

\begin{document}
\firstheader

\noindent\makebox[\paperwidth][l]{%
\begin{sidebarfirst}
  \begin{center}
    \profilephoto{perfil.png}{https://portafolio-tech-eight.vercel.app/}\\[1.2mm]
    {\color{cvmuted}Perfil profesional}
  \end{center}

  \sideheading{Sobre mí}
  Desarrollador full-stack e ingeniero electrónico con experiencia en productos web y móviles, automatización bancaria, infraestructura cloud e IA aplicada. Construyo soluciones end-to-end con Python, FastAPI, React, Next.js y React Native, desde la integración de sistemas hasta su puesta en producción.

  \sideheading{Contacto}
  \contactline{\faEnvelope}{\href{mailto:jguevaral@uni.pe}{jguevaral@uni.pe}}
  \contactline{\faPhone}{+51 991 004 126}
  \contactline{\faMapMarker}{Lima, Perú}
  \contactline{\faLinkedin}{\href{https://www.linkedin.com/in/jhunior-guevara-889483162}{linkedin.com/in/jhunior-guevara-889483162}}
  \contactline{\faGithub}{\href{https://github.com/eduardo202020}{github.com/eduardo202020}}
  \contactline{\faGlobeAmericas}{\href{https://portafolio-tech-eight.vercel.app/}{portafolio-tech-eight.vercel.app}}

  \sideheading{Idiomas}
  \langline{Spain.png}{Español - Nativo}
  \langline{England.png}{Inglés - Intermedio}

  \sideheading{Portafolio web}
  \begin{center}
    \href{https://portafolio-tech-eight.vercel.app/}{\includegraphics[width=3.45cm]{qr-portfolio.png}}\\[1mm]
    \footnotesize Escanea para ver demos, repositorios y proyectos en producción.
  \end{center}
\end{sidebarfirst}
\hfill
\begin{mainfirst}
  \cvsection{\faBriefcase\ EXPERIENCIA}
  \cventry{2026}{Automatización bancaria e infraestructura cloud}{\href{https://careers-inc.nttdata.com/}{NTT DATA}}{Desarrollo y automatización de procesos bancarios con apoyo de IA mediante OneDrive, Teams, Power Automate, macros y Mainframe 3270. Gestión de infraestructura en AWS y Google Cloud con Terraform. En AWS, uso de EC2, S3, IAM, redes VPC/VPN, servicios serverless y bases de datos SQL/MariaDB.}{}
  \cventry{\mbox{2024--2026}}{Full-Stack Software Developer}{Tech Labs - Universidad Nacional de Ingeniería}{Construcción end-to-end de una plataforma en producción con Next.js, React y TypeScript. Implementación de roles, inventario, préstamos, calendario, PostgreSQL, Supabase y visualizaciones 3D.}{Lima, Perú}
  \cventry{\mbox{2022--2024}}{Desarrollador Full-Stack Freelance}{Proyectos propios y por encargo}{Desarrollo de aplicaciones web y móviles, APIs con FastAPI y automatizaciones con Python y LLMs, desde la planificación hasta la implementación y revisión.}{Lima, Perú}
  \cventry{\mbox{2014--2022}}{Ingeniero Electrónico}{Universidad Nacional de Ingeniería}{Investigación y desarrollo de soluciones tecnológicas con Raspberry Pi, sensores, IoT, integración de servicios y herramientas de IA aplicada.}{Lima, Perú}

  \cvsection{\faRocket\ PROYECTOS DESTACADOS}
  \projectentry{Papeleta Clara}{React Native + FastAPI}{Primer lugar en la Hackatón Nacional Transformagob 2026, organizada por la PCM: app para el SAT que facilita consultar papeletas, plazos y evidencias, con orientación por voz o texto.}{https://github.com/eduardo202020/satApp}
  \projectentry{Tech Labs}{Next.js + TypeScript + Supabase}{Plataforma en producción para inventario, préstamos, autenticación, calendario y visualización de recursos tecnológicos.}{https://github.com/eduardo202020/tech-lab}
  \projectentry{MuseIQ + MuseRAG}{React Native + FastAPI + Chroma}{Recorridos de museo con beacons BLE, voz y chat contextual; backend RAG adaptado al museo, sala y obra activa.}{https://github.com/eduardo202020/museApp}
  \projectentry{Gran Chef App}{Aplicación móvil}{Experiencia culinaria con recetas, filtros, favoritos, modo cocina y navegación visual centrada en contenido.}{https://github.com/eduardo202020/food_app_v2}

  \cvsection{\faGraduationCap\ EDUCACIÓN}
  \cventry{}{Ingeniería Electrónica}{Universidad Nacional de Ingeniería}{Formación en electrónica, sistemas embebidos, control, telecomunicaciones, software, automatización e IoT.}{Lima, Perú}
  \cventry{}{Desarrollo de software}{Platzi, Udemy y formación autodidacta}{Aprendizaje continuo en full-stack, aplicaciones móviles, FastAPI, IA, AWS, Terraform, Linux y despliegue de aplicaciones.}{Online}
\end{mainfirst}
}

\newpage
\nextheader
\noindent\makebox[\paperwidth][l]{%
\begin{sidebarnext}
  \sideheading{Fortalezas}
  \begin{center}
    \cvtag{Arquitectura de producto}\cvtag{Automatización}\cvtag{Resolución de problemas}\cvtag{Integración de IA}\cvtag{Prototipado rápido}\cvtag{Trabajo autónomo}
  \end{center}

  \sideheading{Stack principal}
  \begin{center}
    \cvtag{Python}\cvtag{FastAPI}\cvtag{React}\cvtag{Next.js}\cvtag{TypeScript}\cvtag{React Native}\cvtag{Expo}\cvtag{PostgreSQL}\cvtag{MariaDB}\cvtag{Supabase}\cvtag{AWS}\cvtag{Google Cloud}\cvtag{Terraform}\cvtag{Linux Servers}\cvtag{Power Automate}\cvtag{Microsoft 365}\cvtag{Mainframe 3270}\cvtag{LangChain}\cvtag{RAG}\cvtag{LLMs}\cvtag{Raspberry Pi}\cvtag{IoT}
  \end{center}

  \sideheading{Áreas de enfoque}
  \begin{itemize}
    \item IA aplicada
    \item Producto digital
    \item Automatización
    \item Desarrollo web y móvil
    \item RAG
    \item IoT y sistemas embebidos
  \end{itemize}

  \sideheading{Enlaces}
  \begin{itemize}
    \item \href{https://portafolio-tech-eight.vercel.app/}{Portafolio web}
    \item \href{https://github.com/eduardo202020}{GitHub}
    \item \href{https://www.linkedin.com/in/jhunior-guevara-889483162}{LinkedIn}
    \item \href{https://github.com/eduardo202020/satApp}{Papeleta Clara}
    \item \href{https://github.com/eduardo202020/museApp}{MuseIQ}
    \item \href{https://github.com/eduardo202020/tech-lab}{Tech Labs}
    \item \href{https://github.com/eduardo202020/food_app_v2}{Gran Chef App}
  \end{itemize}
\end{sidebarnext}
\hfill
\begin{mainnext}
  \cvsection{\faDesktop\ HABILIDADES TÉCNICAS}
  \skillrow{Frontend y móvil}{\textbf{React, Next.js, TypeScript:} interfaces y productos web.\par\vspace{0.7mm}\textbf{React Native, Expo:} apps móviles, sensores, BLE, voz y navegación.}
  \skillrow{Backend y datos}{\textbf{Python, FastAPI:} APIs, automatización, orquestación y servicios.\par\vspace{0.7mm}\textbf{PostgreSQL, MariaDB, Supabase, SQLite y Chroma:} bases de datos SQL, persistencia, tiempo real y búsqueda vectorial.}
  \skillrow{IA y RAG}{\textbf{LangChain, LLMs, prompts y Chroma:} asistentes contextuales, recuperación de información y apoyo al desarrollo de software.}
  \skillrow{Automatización y cloud}{\textbf{Power Automate, Microsoft 365, macros y Mainframe 3270:} procesos bancarios con apoyo de IA.\par\vspace{0.7mm}\textbf{AWS y Terraform:} despliegue e infraestructura como código con EC2, S3, IAM, redes VPC/VPN, servicios serverless y bases de datos SQL/MariaDB.\par\vspace{0.7mm}\textbf{Linux servers:} entornos de desarrollo y despliegue de aplicaciones.}
  \skillrow{Hardware e IoT}{\textbf{Raspberry Pi, sensores, BLE e IoT:} prototipado, integración hardware-software y captura de datos contextuales.}

  \cvsection{\faImages\ APPS CONSTRUIDAS}
  \apppreview{Papeleta Clara}{previews/7.jpg}{previews/8.jpg}{previews/9.jpg}{Primer lugar en la Hackatón Nacional Transformagob 2026, organizada por la PCM. App para el SAT que permite consultar papeletas, plazos y evidencias, con orientación por voz o texto.}{https://github.com/eduardo202020/satApp}
  \apppreview{MuseIQ}{previews/1.jpg}{previews/2.jpg}{previews/3.jpg}{Recorridos de museo con reconocimiento de sala, escaneo QR, piezas 3D, voz y chat contextual. React Native, Expo, BLE y RAG.}{https://github.com/eduardo202020/museApp}
  \apppreview{Gran Chef App}{previews/4.jpg}{previews/5.jpg}{previews/6.jpg}{Experiencia móvil culinaria con recetas, filtros, favoritos, modo cocina y navegación visual centrada en contenido.}{https://github.com/eduardo202020/food_app_v2}
  \webpreview{Tech Labs}{tech-lab.png}{Plataforma web real para inventario, préstamos, autenticación, calendario y visualización de recursos tecnológicos.}{https://github.com/eduardo202020/tech-lab}

  \cvsection{\faCertificate\ CURSOS Y FORMACIÓN}
  \begin{itemize}
    \item \textbf{LLM Engineering - Udemy.} Construcción de agentes inteligentes con LangChain y herramientas modernas de IA.
    \item \textbf{Platzi y práctica personal:} formación continua en desarrollo full-stack, AWS, Terraform y Linux; uso de IA para documentación, aprendizaje, resolución de incidencias y despliegue de aplicaciones.
  \end{itemize}

\end{mainnext}
}
\end{document}
